\scanpage{105}

Notons qu'on a sur le triple

\[
\left\{\begin{aligned}
 &P=(S,A,F)\qquad\text{des relations ``d'\unsure{incidences|incidence}''}\\
 &R_1\subset S\times A\qquad R=\{(s,a)\in S\times A\mid s\in a\}\\
 &R_2\subset A\times F\qquad R_2=\{(a,f)\in A\unsure{\times|-} F\mid a\subset f\}\\
 &R_3=R_2\circ R_1\subset S\times F\qquad R_3=\{(s,f)\in S\times F\mid s\in f\},
\end{aligned}\right.\tag*{(10)}
\]

% slanted
\margin{\uline{NB} $R'_2$ de (5) est la relation symétrique de $R_2$ de (10)}

qui satisfont ``l'axiome numérique''\ des structures 2-polyédrales

\[
\left\{\begin{aligned}
 \text{a)}\ &\forall a\in A,\ \text{il existe exactement deux sommets incidents}\\
 &\text{et deux faces incidentes}\\
 \text{b)}\ &\forall(s,f)\in R_3,\ \text{il existe exactement deux arêtes}\\
 &\text{incidentes à la fois à }s\text{ et à }f
\end{aligned}\right.\tag*{(11)}
\]

et aussi l'axiome de connexité en les faces -- qui sont en effet des triangles.

Considérons d'autre part la question de connexité en \deleted{le}\inserted{un} sommet $s\in S$, en introduisant son ``étoile''\ $\operatorname{Et}(s)$ ou $A(s)$

\[
\deleted{\operatorname{Et}}A(s)=\{t\in S\mid\{s,t\}\in A\},\tag*{(12)}
\]

qui est \deleted{muni}\inserted{l'ens des sommets} d'un \deleted{structure} \uline{contour combinatoire}\inserted{$\deleted{\unsure{A(s)}}P(s)$}, dont l'ens des arêtes \deleted{est}\inserted{peut se décrire comme}

% slanted
\margin{$P(s)$ peut être considéré comme le localisé de la structure \unsure{triangul. |triangulaire} $P$ en $s$}

\[
\begin{aligned}
 &\deleted{\unsure{A}} F(s)=\bigl\{\{t,t'\}\in\unsure{\mathscr P|\mathcal P|\mathfrak P}_2(\deleted{\unsure{F}}S)\deleted{\unsure{(s)}}\mid\{s,t\unsure{.|,}t'\}\in F\bigr\}\\
 &\qquad\rotatebox{90}{$\simeq$}\\
 &\deleted{\unsure{f_s\in}}R_3(s)=\{f\in F\mid s\in f\}.
\end{aligned}\tag*{(13)}
\]

On a bien sur \deleted{des} décompositions
% archive sheet number: 104
